\documentclass[11pt]{article}
\usepackage{fltcourse}

\title{Level 8: Automorphic forms and potential automorphy}
\author{Kevin Buzzard}
\date{}

\begin{document}
\maketitle

\noindent
We have reduced Fermat's Last Theorem to $B_7=B_{6a}\wedge B_{6b}$: an
irreducible hardly ramified mod $\ell$ representation lifts to an $\ell$-adic
one ($B_{6a}$), and, once its residual representation descends to $\F_\ell$,
spreads out into a compatible family ($B_{6b}$). Both are theorems about Galois
representations with no modularity hypothesis, and both are out of reach by
purely Galois-theoretic means. The way forward is to bring in the
\emph{automorphic} world, where the analogues of lifting and spreading out are
easy, and to connect the two worlds.

\medskip\noindent
This level introduces just enough of the automorphic side to state the three
statements $B_{8a}$, $B_{8b}$, $B_{8c}$ that make up $B_8$, and proves that
$B_8$ implies $B_7$ (and hence FLT). As in Level 5, the definitions on the
automorphic side are standard --- quaternion algebras, ad\`eles, automorphic
forms, Hecke operators are all part of the literature (and of the formalized
library) --- so we recall them briefly rather than build them up. The genuinely
new mathematical content of the level is a theorem, proved in full, on the
absolute irreducibility of hardly ramified representations.

\section{Automorphic forms on a definite quaternion algebra}

Let $F$ be a totally real number field (every embedding $F\hookrightarrow\C$
lands in $\R$) with ring of integers $A$, and let $D$ be a \emph{totally
definite} quaternion algebra over $F$: a four-dimensional central division
algebra with $D\otimes_{F,v}\R\cong\bbH$ (Hamilton's quaternions) at every real
place $v$. Let $\adele{F}=\prod'_P F_P$ be the finite ad\`eles of $F$
(restricted product over the finite places $P$, with respect to the integer
rings $A_P$), and put $D_f:=D\otimes_F\adele{F}$, a topological ring containing
$D$. Its unit group $D_f^\times$ is a locally compact topological group
containing $D^\times$.

Total definiteness makes the archimedean symmetric space a point, so there is
no analysis: for a coefficient ring or field $M$, a (weight-two, trivial
coefficient) \emph{automorphic form} for $D$ valued in $M$ is simply a locally
constant function $\phi\colon D_f^\times\to M$ with $\phi(\delta g)=\phi(g)$ for
all $\delta\in D^\times$; equivalently a locally constant function
$D^\times\backslash D_f^\times\to M$. Write $\calA_D(M)$ for the space of these;
$D_f^\times$ acts by right translation, $(x\bullet\phi)(g)=\phi(gx)$. For an
open compact subgroup $U\subseteq D_f^\times$ (a \emph{level}) put
$\calA_D(U;M):=\calA_D(M)^U$, the forms fixed by $U$. The essential finiteness
input (Fujisaki; known in the 1960s) is that $D^\times\backslash D_f^\times$ is
compact, so $D^\times\backslash D_f^\times/U$ is finite and
$\calA_D(U;M)\cong M^{\,n}$ with $n=|D^\times\backslash D_f^\times/U|$; when $M$
is a commutative ring this is a finite free $M$-module.

\emph{Hecke operators.} For $g\in D_f^\times$ and compact open subgroups $U,V$,
the double coset $UgV=\coprod_{i=1}^{m} g_iV$ is a finite union of left cosets,
and $[UgV]\phi:=\sum_i g_i\bullet\phi$ defines an $M$-linear map
$\calA_D(V;M)\to\calA_D(U;M)$. Fix a maximal order in $D$ and, for each finite
place $P$ splitting $D$, an identification with $\GL_2$ locally; this gives, for
each finite set $S$ of finite places splitting $D$, the level
$U_0(S)$ (matrices upper-triangular mod $P$ for $P\in S$), and for each finite
place $P\notin S$ splitting $D$ the operator
$T_P:=[U_0(S)\,\pi_P\,U_0(S)]$, where $\pi_P$ is the id\`ele
$\bigl(\begin{smallmatrix}\varpi_P&0\\0&1\end{smallmatrix}\bigr)$ at $P$ and $1$
elsewhere. The $T_P$ commute. A \emph{system of automorphic $M$-eigenvalues for
$(F,D,U_0(S))$} is the family $(t_P)_{P\notin S}$ of eigenvalues of a common
eigenvector $\phi\in\calA_D(U_0(S);M)$ (an \emph{eigenform}): $T_P\phi=t_P\phi$.

\begin{example}
The constant function $\phi\equiv1$ is an eigenform of every level, with
$T_P\phi=(1+|A/P|)\phi$; so $t_P=1+|A/P|$. If $F=\Q$ and $P=(p)$ this is
$t_p=1+p$, the trace of $\Frob_p$ on $\mathbf1\oplus\chi$.
\end{example}

The bridge to Galois representations is the following deep theorem, whose proof
--- via the Jacquet--Langlands correspondence, the trace formula, and the
cohomology of Shimura varieties --- was completed by Taylor in 1989. It is the
single most substantial classical input to the whole proof, and we assume it.

\begin{theorem}[Taylor, 1989]\label{thm:taylor}
Let $k$ be a finite field of characteristic $\ell$, or a finite extension of
$\Ql$. Let $F$ be totally real, $D/F$ totally definite ramified at the finite
set $S_0$ of finite places, and let $(t_P)$ be a system of automorphic
$k$-eigenvalues for $(F,D,U_0(S))$ with $S\cap S_0=\varnothing$. Then there is a
continuous semisimple representation $\rho_t\colon\GF\to\GL_2(k)$ such that
$\rho_t$ is unramified outside $S\cup S_0\cup\{P:\ell\in P\}$, and
$\tr\rho_t(\Frob_P)=t_P$ for every $P$ outside this set. Its determinant is the
cyclotomic character up to a finite-order character.
\end{theorem}

A representation $\GF\to\GL_2(k)$ arising this way is called \emph{automorphic}.

\section{Potential automorphy}

We cannot expect a mod $\ell$ representation of $\GQ$ to be automorphic on the
nose; but it may become so after restriction to a suitable totally real field.
Let $k$ be a finite field of characteristic $\ell$ and $\rhobar\colon\GQ\to\GL_2(k)$
continuous, and let $K_{\rhobar}$ be the finite Galois extension of $\Q$ cut out
by $\ker\rhobar$.

\begin{definition}\label{def:potaut}
We say $\rhobar$ is \emph{potentially automorphic} if there exist
\begin{itemize}
\item a totally real field $F$, Galois over $\Q$, unramified at $\ell$, and
\emph{disjoint} from $K_{\rhobar}$ \textup{(}$F\cap K_{\rhobar}=\Q$\textup{)};
\item a totally definite quaternion algebra $D/F$ split at all finite places
\textup{(}so $[F:\Q]$ is even\textup{)};
\item a finite set $S$ of finite places of $F$ coprime to $\ell$, and a system
$(t_P)$ of automorphic $k$-eigenvalues for $(F,D,U_0(S))$;
\end{itemize}
such that the restriction $\rhobar|_{\GF}$ satisfies
$\tr\rhobar|_{\GF}(\Frob_P)=t_P$ for every finite place $P\notin S$ with
$\ell\notin P$.
\end{definition}

The disjointness $F\cap K_{\rhobar}=\Q$ ensures that $\rhobar|_{\GF}$ has the
same image as $\rhobar$, so that no information is lost on restriction.

\section{The three heads}

Set $B_8:=B_{8a}\wedge B_{8b}\wedge B_{8c}$. In all three, $\ell\geq5$.

\begin{Bn}{B_{8a}}
\textup{(Automorphic lifting.)} Let $k$ be finite of characteristic $\ell$ and
let $\rhobar$ be a hardly ramified $k$-representation which is potentially
automorphic and such that $\rhobar|_{\Gal(\Qbar/\Q(\zeta_\ell))}$ is absolutely
irreducible. Then $\rhobar$ lifts to a hardly ramified $\ell$-adic
representation.
\end{Bn}

\begin{Bn}{B_{8b}}
\textup{(Automorphic spreading out.)} Let $R$ be the ring of integers of a
finite extension of $\Ql$, with residue field $k$, and let $\rho$ be a hardly
ramified $R$-representation whose reduction $\rhobar$ is potentially automorphic
and has $\rhobar|_{\Gal(\Qbar/\Q(\zeta_\ell))}$ absolutely irreducible. Then
$\rho$ spreads out into a compatible family of hardly ramified $p$-adic
representations.
\end{Bn}

\begin{Bn}{B_{8c}}
\textup{(Potential automorphy.)} Every irreducible hardly ramified
$\Z/\ell\Z$-representation of $\GQ$ is potentially automorphic.
\end{Bn}

\noindent
The changes from $B_{6a},B_{6b}$ are exactly the replacement of ``residual
representation defined over $\F_\ell$'' by the pair of conditions ``potentially
automorphic'' and ``absolutely irreducible on $\Q(\zeta_\ell)$''. To recover
$B_{6a},B_{6b}$ we must know that an irreducible hardly ramified
$\Z/\ell\Z$-representation satisfies both.

\section{The absolute irreducibility theorem}

This is the one part of the level proved from scratch. It uses one piece of
local arithmetic-geometry --- a bound of Raynaud, from his 1974 paper on group
schemes of type $(p,\dots,p)$ --- which we assume; everything else is group
theory and linear algebra.

Throughout, $k$ is a finite field of characteristic $\ell\geq5$ with algebraic
closure $\kbar$, and $\chibar\colon\GQ\to\Fl^\times\subseteq\kbar^\times$ is the
mod $\ell$ cyclotomic character; put $G=\GQ$ and $H=\Gal(\Qbar/\Q(\zeta_\ell))$,
so $G/H\cong\Fl^\times$ is cyclic of order $\ell-1$. Fix an embedding
$\Qbar\hookrightarrow\Qlbar$, giving a decomposition group $D=\GQl\subseteq G$
with inertia subgroup $I$.

\subsection*{Assumed local inputs}
We assume the following, all of which were available in the 1970s.
\begin{itemize}
\item[(L1)] \emph{(Raynaud's bound.)} Let $K/\Ql$ be a \emph{ramified quadratic}
extension and $\psi\colon\Gal(\Qlbar/K)\to\kbar^\times$ a flat character. Then
there is an integer $d$ with $0\leq d\leq2$ such that
$\psi(\sigma^2)=\chibar(\sigma)^d$ for all $\sigma\in I$. \textup{(}Here
$e(K/\Ql)=2<\ell-1$, the range where Raynaud's bound has content.\textup{)}
\item[(L2)] \emph{(Flatness of stable-line characters.)} If $\rhobar$ is flat
at $\ell$ and $K/\Ql$ is finite, any character by which $\Gal(\Qlbar/K)$ acts on
a stable $\kbar$-line is a flat character of $\Gal(\Qlbar/K)$.
\item[(L3)] \emph{(Complex conjugation.)} There is $c\in G$ with $c^2=1$ and
$\chibar(c)=-1$.
\item[(L4)] \emph{(Total ramification.)} $\chibar(I)=\Fl^\times$.
\end{itemize}

\begin{theorem}\label{thm:absirr}
Let $\rhobar\colon\GQ\to\GL_2(k)$ be continuous, irreducible, flat at $\ell$,
with $\det\rhobar=\chibar$ \textup{(}for instance, any irreducible hardly
ramified $k$-representation with $\ell\geq5$\textup{)}. Then
$\rhobar|_{H}=\rhobar|_{\Gal(\Qbar/\Q(\zeta_\ell))}$ is absolutely irreducible.
\end{theorem}

\begin{proof}
Suppose not. We work over $\kbar$, on $\overline V=\kbar^2$.

\emph{Step 1: $\rhobar$ is absolutely irreducible.} As in Level 6, using
$c$ from (L3): $\rhobar(c)$ is an involution of determinant $\chibar(c)=-1$,
with distinct eigenvalues $\pm1$ (as $\ell>2$) whose eigenlines are defined
over $k$. Any $\rhobar(G)$-stable $\kbar$-line would be $\rhobar(c)$-stable,
hence one of these $k$-rational lines, contradicting irreducibility over $k$.
So $\rhobar\otimes\kbar$ is irreducible.

\emph{Step 2: Clifford theory.} $H\trianglelefteq G$ with $G/H$ cyclic, and by
hypothesis $(\rhobar\otimes\kbar)|_H$ is reducible. An $H$-stable line
$\Lambda_1$ has $G$-translates that are again $H$-stable ($H$ normal); their sum
is $G$-stable, hence all of $\overline V$ by Step 1, so
$\overline V=\Lambda_1\oplus\Lambda_2$ with $H$ acting on $\Lambda_i$ by a
character $\psi_i$. If $\psi_1=\psi_2$ then $H$ acts by scalars, every line is
$H$-stable, and taking $g\in G$ lifting a generator of $G/H$, an eigenline of
$\rhobar(g)$ is $\{g,H\}$-stable, hence $G$-stable: contradiction. So
$\psi_1\neq\psi_2$.

\emph{Step 3: the index-two subgroup.} Since $\psi_1\neq\psi_2$, the only
$H$-stable lines are $\Lambda_1,\Lambda_2$, and $G$ permutes them, nontrivially
by Step 1. Thus $L:=\{g:g\Lambda_1=\Lambda_1\}$ has index $2$ in $G$, contains
$H$, acts on $\Lambda_1$ by a character $\psi$ extending $\psi_1$, and every
$g\notin L$ swaps $\Lambda_1,\Lambda_2$. As $G/H\cong\Fl^\times$ is cyclic of
even order $\ell-1$, its unique index-two subgroup is the squares, so
$L=\chibar^{-1}\!\bigl((\Fl^\times)^2\bigr)$.

\emph{Step 4: a good inertia element.} By (L4) pick $\tau\in I$ with
$\chibar(\tau)$ a generator of $\Fl^\times$. Then $\chibar(\tau)$ is a nonsquare,
so $\tau\notin L$; it has order $\ell-1\geq4$, so $\chibar(\tau)\neq-1$; and
$\tau^2\in L$.

\emph{Step 5: the determinant at $\tau$.} Take $0\neq v\in\Lambda_1$ and
$w:=\rhobar(\tau)v$, which spans $\Lambda_2$ (as $\tau$ swaps the lines). Then
$\rhobar(\tau)w=\rhobar(\tau^2)v=\psi(\tau^2)v$, so in the basis $(v,w)$,
$\rhobar(\tau)=\left(\begin{smallmatrix}0&\psi(\tau^2)\\1&0\end{smallmatrix}\right)$
and $\det\rhobar(\tau)=-\psi(\tau^2)$. Since $\det\rhobar=\chibar$,
\[
  \chibar(\tau)=-\psi(\tau^2). \qquad(*)
\]

\emph{Step 6: the local bound.} As $\tau\in D\setminus L$, the group
$D\cap L=\Gal(\Qlbar/K)$ for a quadratic $K/\Ql$; and $\tau\in I\setminus L$
shows $K/\Ql$ is \emph{ramified}. The line $\Lambda_1$ is stable under
$\rhobar(D\cap L)$ via $\psi|_{D\cap L}$, which is flat by (L2). By (L1) there is
$d\in\{0,1,2\}$ with $\psi(\sigma^2)=\chibar(\sigma)^d$ for all $\sigma\in I$; in
particular $\psi(\tau^2)=\chibar(\tau)^d$.

\emph{Step 7: contradiction.} Write $x=\chibar(\tau)$, of order $\ell-1\geq4$.
By $(*)$, $x=-x^{d}$. If $d=1$ then $1=-1$, impossible as $\ell$ is odd. If
$d=0$ then $x=-1$, of order $2$, contradicting order $\geq4$. If $d=2$ then
$x=-x^2$, so $1=-x$, i.e.\ $x=-1$, the same contradiction. Hence $\rhobar|_H$ is
absolutely irreducible.
\end{proof}

\begin{remark}
The theorem genuinely fails for $\ell=3$: there are irreducible flat mod $3$
representations with cyclotomic determinant induced from $\Q(\sqrt{-3})=\Q(\zeta_3)$
whose restriction to $\Gal(\Qbar/\Q(\zeta_3))$ is reducible. This is why the
prime $3$ was handled separately, and why $B_{8a},B_{8b}$ require $\ell\geq5$.
\end{remark}

\section{The boss: $B_8$ implies FLT}

\begin{theorem}[Boss of Level 8]\label{thm:boss8}
Statements $B_{8a}$, $B_{8b}$ and $B_{8c}$ together imply Fermat's Last
Theorem, modulo results known in the 1980s.
\end{theorem}

\begin{proof}
By the boss of the previous level it suffices to prove $B_{6a}$ and $B_{6b}$.

For $B_{6a}$: let $\rhobar$ be an irreducible hardly ramified
$\Z/\ell\Z$-representation, $\ell\geq5$. By $B_{8c}$ it is potentially
automorphic, and by Theorem~\ref{thm:absirr} (which applies, a hardly ramified
representation being flat at $\ell$ with cyclotomic determinant)
$\rhobar|_{\Q(\zeta_\ell)}$ is absolutely irreducible. Hence $B_{8a}$ applies
and $\rhobar$ lifts to a hardly ramified $\ell$-adic representation: this is
$B_{6a}$.

For $B_{6b}$: let $\rho$ be a hardly ramified $\calO$-representation whose
reduction $\rhobar$ is the base change of an irreducible hardly ramified
$\Z/\ell\Z$-representation $\rhobar_0$. By $B_{8c}$, $\rhobar_0$ is potentially
automorphic; potential automorphy depends only on the traces of Frobenius,
which are unchanged under the base change $\rhobar_0\rightsquigarrow\rhobar$, so
$\rhobar$ is potentially automorphic as well. By Theorem~\ref{thm:absirr},
$\rhobar_0|_{\Q(\zeta_\ell)}$ --- and hence its base change
$\rhobar|_{\Q(\zeta_\ell)}$ --- is absolutely irreducible. So $B_{8b}$ applies
and $\rho$ spreads out into a compatible family: this is $B_{6b}$.

Therefore $B_8\Rightarrow B_7\Rightarrow{}$FLT.
\end{proof}

\begin{remark}
Of the three heads, $B_{8a}$ and $B_{8b}$ are, on the automorphic side, soft:
lifting a system of Hecke eigenvalues from characteristic $\ell$ to
characteristic zero is the going-down theorem, and spreading a
characteristic-zero system to another prime is elementary linear algebra over a
number field. The difficulty is entirely in the word ``potentially'': these
easy automorphic facts only control $\rho|_{\GF}$, not $\rho$ itself. Turning
control of $\rho|_{\GF}$ into control of $\rho$ requires an automorphy lifting
theorem, which is the subject of the remaining levels; $B_{8c}$ needs it too.
\end{remark}

\end{document}
